\section{Spiritual waterboarding}

For Julius Evola, philosophy was never a dry academic discipline, a mincing of words, a defense and attack of various doctrines, or trite rhetoric. For him, it was a matter of self-realisation that transcended the categories of thought, and this by necessity stands above any logic or rationality.

In his \emph{Saggi sull Idealismo Magico} \emph{Essays on Magical Idealism} , Evola relates this Hindu story:

\begin{quotex}
A disciple asked his spiritual teacher, while they were bathing in the river, when he would finally be able to realise \emph{Brahman}. The teacher, instead of answering, pushed his head under water and held him down until, feeling himself drowning, the student freed himself and re-emerged. The teacher then explained: ``When the desire in you to realise \emph{Brahman} becomes as intense and deep as how you were just driven to reassert your physical life---only then will you achieve satisfaction.''
\end{quotex}


Evola goes on to explain:

\begin{quotex}
That expresses the whole of magical development: as long as the will and the desire to realise oneself remains in mental shadows, as long as it does not equal, by penetrating into all one's being, the intensity of that obscure and irrational power that is asserted at the core of our organism, there is no point in expecting any concrete result. With some justice, Jacobi noted that man cannot improve himself by means of ideas or reasoning; what he needs is to be organised and therefore to organise himself.
\end{quotex}


This explains why Evola is often so misunderstood. He is not teaching a new doctrine or recommending a specific religion. Rather, he is trying to awaken in the reader the desire for transcendence, which can only arise in the deepest parts of himself, prior to ideas, rationality, or convention. Without making the interior efforts at such transcendence, Evola's entire scheme will remain opaque.

What follows is the opening paragraph to chapter V of the essays:

\paragraph{The Essence of Magical Development}


\begin{quotex}
I said, ``You are gods''
\flright{\textsc{John 10:34}}

What distinguishes magical idealism is its essentially practical character: its fundamental requirement is not to substitute one intellectual conception of the world with another, but rather to create in the individual a new ``dimension'' and a new depth of life. Certainly, it does turn into an abstract opposition between theory and practice; already in the theoretical and cognitive as such---and therefore in that which only be given to be revealed to a reader--- it sees a level of creative activity, therefore it holds that such a level represents only an outline, a beginning of an effort in respect to a phase of deeper realisation, which is that of magic or practice properly called, in which the former must be continued and fulfilled. The lowering of the ``being'' of ontology and gnoseology to a ``having to be'' and the development of the activity of the judgment of value, into which the same theoretical judgment is accordingly transformed, until the judgment of existence, to a cosmically creative act of faith, such is the essence of the present doctrine. Whoever therefore does not know how to convert the principles of magical idealism into power agents in his interior being, to profound needs that impel him towards a concrete and living realisation, he kills that idealism by the tritest rhetoric. The one who can truly say to himself: ``If magic didn't exist, today I must create it by myself'', is rather welcome here.
\end{quotex}



\flrightit{Posted on 2007-12-20 by Cologero}
